% --------------------------------------------------------------
% This is all preamble stuff that you don't have to worry about. Head down to
% where it says "Start here"
% --------------------------------------------------------------
 
\documentclass[12pt]{article}
 
\usepackage[margin=1in]{geometry} 
\usepackage{amsmath,amsthm,amssymb}
\usepackage{hyperref}
\usepackage{listings}
\usepackage{xcolor}

\newcommand{\todobox}[3]{%
	\colorbox{#1}{\textcolor{white}{\sffamily\bfseries\scriptsize #2}}%
	~\textcolor{#1}{#3} %
	\textcolor{#1}{$\triangleleft$}%
} \newcommand{\flavio}[1]{\todobox{violet}{Fla:}{#1}}

\lstdefinelanguage{bash}{ morekeywords={sudo, bash, echo, cd, ls, rm, cp, mv},
  sensitive=true, morecomment=[l]{\#}, morestring=[b]", }

\lstset{ language=bash, basicstyle=\ttfamily\small,
  backgroundcolor=\color{gray!8}, frame=none, breaklines=true,
  postbreak=\mbox{\textcolor{red}{$\hookrightarrow$}\space}, }
 
\newcommand{\N}{\mathbb{N}} \newcommand{\Z}{\mathbb{Z}}
 
\newenvironment{theorem}[2][Theorem]{\begin{trivlist}
\item[\hskip \labelsep {\bfseries #1}\hskip \labelsep {\bfseries
#2.}]}{\end{trivlist}}
\newenvironment{lemma}[2][Lemma]{\begin{trivlist}
\item[\hskip \labelsep {\bfseries #1}\hskip \labelsep {\bfseries
#2.}]}{\end{trivlist}}
\newenvironment{exercise}[2][Exercise]{\begin{trivlist}
\item[\hskip \labelsep {\bfseries #1}\hskip \labelsep {\bfseries
#2.}]}{\end{trivlist}}
\newenvironment{problem}[2][Problem]{\begin{trivlist}
\item[\hskip \labelsep {\bfseries #1}\hskip \labelsep {\bfseries
#2.}]}{\end{trivlist}}
\newenvironment{question}[2][Question]{\begin{trivlist}
\item[\hskip \labelsep {\bfseries #1}\hskip \labelsep {\bfseries
#2.}]}{\end{trivlist}}
\newenvironment{corollary}[2][Corollary]{\begin{trivlist}
\item[\hskip \labelsep {\bfseries #1}\hskip \labelsep {\bfseries
#2.}]}{\end{trivlist}}

\newenvironment{solution}{\begin{proof}[Solution]}{\end{proof}}
 
\begin{document}
 
% --------------------------------------------------------------
%                         Start here
% --------------------------------------------------------------
 
\title{Assignment Sheet 1}
\author{Flavio Toffalini\\ %replace with your name
Advanced Automatic Testing (211067-SoSe 2026)}

\maketitle

\noindent This is the first Assignment Sheet for the course Advanced Automatic
Testing (211067-SoSe 2026). The assignment will be available for a week. Please
follow the submission instructions carefully or contact the instructor for any
question.

\paragraph{Summary:}
\begin{itemize}
    \item \textbf{Uptime:} From 22th of April 2026 at 23:59 to 30th of April 2026 at 23:59.
    \item \textbf{Instructor:} Flavio Toffalini (flavio.toffalini@rub.de).
\end{itemize}





% \begin{enumerate} %the first video is required (but you don't have to watch it
% first), you can choose the other 4 from the list on the Course Materials page
% \item Grit: the power of passion and perseverance by Angela Lee Duckworth  
% \item Replace with title for video 2 \item Replace with title for video 3
% \item Replace with title for video 4 \item Replace with title for video 5
% \end{enumerate}

% Replace this text with your reflection.  Your total reflection should be at
% least 15 sentences long. You should comment on each of the videos you watched,
% but rather than reflecting on each video separately, try to reflect on growth
% mindset, productive failure, and grit, in general.

\section*{Remote Environment}

The tasks will be solved in a remote environment. You will be provided with a
Docker container that contains all the necessary tools and libraries to complete
the tasks. The container is hosted on a server, and you will access it via SSH.

\begin{itemize}
    \item \textbf{Register a user}\\
    To register for our remote environment, please visit the following link:\\
    \url{https://moodle.ruhr-uni-bochum.de/mod/data/view.php?id=3617068}\\    
    You need to provide your \emph{matriculation number}, \emph{username}, and
    \emph{public SSH key}. The link also contains instructions on how to
    generate your SSH key. Once registered, please contact the instructors to
    get your user activated. If you have already registered, you can skip this
    step.
    \item \textbf{Connect to a task}\\
    We prepared a server to which you can connect via SSH. The server is
    accessible via RUB VPN. You can find more information in the following
    link:\\
    \url{https://noc.rub.de/web/vpn}\\
    Once your user is activated, you can connect to the server via SSH with the
    following command:\\
    \begin{lstlisting}
ssh -i <path_to_private_keyfile> -p 25222 \
    <username>@server1.ig35s.ruhr-uni-bochum.de
    \end{lstlisting}
    Where \texttt{username} is your username. You will be redirected into a
    Docker container with all the basic tools to solve the tasks.
    \item \textbf{Download Material}\\
    It is possible to transfer files in the remote environment through
    \texttt{scp}. According to your \texttt{scp} version, one of the following
    command should work:
    \begin{lstlisting}
# OpenSSH version >= 9.0        
scp -O -i <path_to_private_keyfile> -P 25222 
    <username>@server1.ig35s.ruhr-uni-bochum.de:/path/file.txt .
file.txt                          100%   27     1.6KB/s   00:00

# OpenSSH version < 9.0        
scp -i <path_to_private_keyfile> -P 25222 
    <username>@server1.ig35s.ruhr-uni-bochum.de:/path/file.txt .
file.txt                          100%   27     1.6KB/s   00:00
    \end{lstlisting}
    Please note that the usage of the \texttt{-O} option \emph{should} work for
    OpenSSH versions equal to or higher than 9.0. You can run \texttt{ssh -V} to
    check your installation. If you're unsure, try both options. If you are in a
    non-Linux environment and can't figure out how to transfer files, please
    contact the instructor.
\end{itemize}

\paragraph{Development:}
All the developed code must be saved in the home folder
\texttt{/home/<username>}, in which the student can create a
\texttt{<submission\_directory>} where storing the produced material. The home
folder is mounted in the host. Therefore, the files stored in this location
remain persistent even after the Docker is destroyed.

\paragraph{Tools Installed:}
The Docker container contains the following tools already installed: AFL++ (in
\texttt{/data/AFLPlusPlus}), gdb, clang, gcc, LLVM, (neo)vim, and other basic
tools. The container is based on Ubuntu 24.04. The user is \texttt{root},
however, the installation of additional tools should be agreed with the
instructor.

 
\section*{Task1: Hello World!}

The first task has the objective to familiarize the students with the
environment and AFL++. 

\begin{itemize}
    \item Write a basic ``Hello, World!'' program in C.
    \item Introduce a bug (e.g., out-of-bounds access) that can be controlled by
    an adversary.
    \item Compile it for AFL++ and fuzz it.
    \item Demonstrate how the bug manifests using \texttt{gdb}.
\end{itemize}

\section*{Task2: Your First Campaign}

The folder \texttt{/data/binutils-2.44} contains the binutils source code. By
starting from that folder, you have to:

\begin{itemize}
    \item Compile the library with AFL++.
    \item Execute a fuzzing campaign on \texttt{objdump}, the campaign will last
    for 1h max.
    \item Produce a coverage report over time and include it in your report.
    \item Find and report interesting crashes
\end{itemize}

\noindent \textbf{Tips:}
\begin{itemize}
    \item Use \texttt{afl-gcc-fast} for compiling
    \item Use \texttt{timeout} for stopping the campaign programmatically
\end{itemize}

\section*{Task3: Bug Hunting 101}

We provide you with a precompiled binary located at
\texttt{/data/libxml2/xmllint}. The program has already been compiled with AFL++
instrumentation. An artificial bug has been injected into \texttt{xmllint}. This
bug is not documented anywhere online. Your objective is to set up a fuzzing
campaign to discover the bug. Approximately one hour of fuzzing should be
sufficient. You may tweak the AFL++ configuration to improve code coverage and
fuzzing efficacy.


% --------------------------------------------------------------
%     You don't have to mess with anything below this line.
% --------------------------------------------------------------
 
\end{document}
